\subsection{The Cartan chart under an actual image hypothesis}
\label{sec:boundary-nonsplit-image-proof}

This subsection constructs the relevant invariant-ring map from
an actual nonsplit-Cartan image and derives its incompatible
ramification divisibility. We first make the coefficient-field
scope of the universal frame map explicit.

\begin{lemma}[Scalar-compatible universal frames]
\label{lem:scalar-compatible-frame}
Put $k_0=\overline{\mathbf F}_p$, and let $R_n(k_0)$ be the universal
full-level deformation ring of a fixed one-dimensional height-two
framing group over $k_0$. Let $A$ be a complete mixed-characteristic
discrete valuation ring with algebraically closed residue
$\kappa\supset k_0$ and a fixed compatible coefficient map
$W(k_0)\subset W(\kappa)\to A$.
A good formal group over $A$ with the scalar-extended special
framing, together with a full $p^n$-frame over a finite extension
valuation ring $B$, gives a continuous homomorphism
\[
                           R_n(k_0)\longrightarrow B.
\]
The level-zero deformation and its identity-reducing formal
isomorphism can be chosen over $A$. If the fraction field of $B$
is the Galois splitting field, the displayed map is equivariant
for the actual Galois image acting on the frame. The integral
formal coordinate used in this evaluation preserves the valuation
of every positive-valuation point.
\end{lemma}

\begin{proof}
The printed deformation category of
\citet[Propositions~4.1--4.3, pp.~570--573]{Drinfeld1974} has residue
field $k_0$. The coefficient extension needed here follows from
the coefficientwise proof; we do not attribute a broader literal
statement to that source.

In the level-zero proof, the difference between successive
deformations over a square-zero layer
$M=\mathfrak m_A^a/\mathfrak m_A^{a+1}$ is a formal-module
$2$-cocycle with coefficients in $M$. Such a layer may be viewed
as an arbitrary $k_0$-vector space. The equations for the cocycle
and its coboundaries have coefficients in the original framing
group over $k_0$. The degree-elimination assertions~(a) and~(b)
in that proof, pp.~571--572, only invert nonzero coefficients of
$k_0$. They eliminate coefficients one degree at a time; each
fixed degree is modified only finitely often. They therefore
apply to $M$-valued coefficients as well as scalar coefficients.
There is no interchange of an infinite product of power-series
coefficients with a tensor product. At height two, one parameter of deformation classes remains,
represented by a change of the single universal parameter.

Starting from the given special framing, choose that parameter
modulo successive powers of $\mathfrak m_A$ to correct the
remaining cocycle class. The same degree elimination gives the
formal coordinate isomorphism at each step. The rigidity argument
of Proposition~4.1 is valid with these coefficients: applying
separating $k_0$-linear functionals to an infinitesimal
homomorphism reduces it to the scalar argument, and finite height
excludes a nonzero homomorphism to the additive formal group.
Thus an isomorphism reducing to the identity is unique. The
successive choices are compatible, and completeness gives the
parameter and formal isomorphism over $A$, hence the map
$R_0(k_0)\to A$. Its linear coordinate coefficient is a unit.
For a positive-valuation argument, the linear term then has
strictly smaller valuation than every higher term, proving
coordinate invariance of the valuation.

The higher-level construction in Proposition~4.3 adjoins roots of
successive primitive division factors at level one and then roots
of $[p](X_i)-x_i$ at later levels. The chosen actual frame supplies
these roots. At a primitive step, its next point is outside the
span already chosen; in the characteristic-zero fraction field,
all corresponding linear factors are nonzero at that point.
Since the full division series vanishes there, the remaining
primitive factor vanishes. The later division relations hold by
the definition of a compatible full frame. The root-adjunction
construction therefore gives the asserted map $R_n(k_0)\to B$.
The level-zero data were constructed over $A$, and their
identity-reducing isomorphism is unique. Change of frame is
functorial. These facts give the stated Galois equivariance.
\end{proof}

\begin{proof}[Proof of Corollary~\ref{cor:boundary-nonsplit-cartan-image}]
Work over the completed maximal unramified base, with algebraically
closed residue field $\kappa$, and let $L$ be the common tame good
extension of degree $m=p^2-1$. The residual representations split
into $\mathbf F_p$-valued characters by
Theorem~\ref{thm:first-boundary-normal-form}. Their wild inertia
is trivial and their tame orders divide $p-1$; hence first torsion
is rational over $L$. We keep the universal deformation and its
Cartan charts over $k_0=\overline{\mathbf F}_p\subset\kappa$.

Put
\[
 c_0=\frac{p-1}{2},\qquad A_0=p^2+p+1,\qquad B_0=p^2.
\]
The normal form gives $z=c_0p^{r+3}$ and
$e=c_0(p^2+1)p^r$. The canonical and noncanonical first-torsion
valuations are therefore
\begin{align*}
 A&=(p+1)e-z=c_0p^r A_0,\\
 B&=z/p=c_0p^r B_0.
\end{align*}
Both models are central. Lemma~\ref{rc-lem:newton} shows that every
exact $p^n$-point above a nonzero canonical, respectively
noncanonical, first-torsion point has valuation $A/p^{2n-2}$,
respectively $B/p^{2n-2}$. Here $v_L(L^\times)=\mathbf Z$ and
$v_L(p)=me$.

Suppose that the finite image for the first group $G$ is contained
in a conjugate of the principal nonsplit Cartan $C$. Write
$R=\mathbf Z/p^n\mathbf Z$ and realize its quadratic algebra as
$R[J]$, where $J^2=\epsilon I$ and the residue of $\epsilon$ is a
nonsquare. Choose a primitive point $v$ whose $p^{n-1}$-multiple
is a nonzero canonical point. The reduction of $J$ has no
$\mathbf F_p$-eigenline. Thus $(v,Jv)$ is an $R$-basis and
$p^{n-1}Jv$ is noncanonical. In this basis
\[
                          J=\begin{pmatrix}0&\epsilon\\1&0\end{pmatrix},
\]
so the Cartan has the standard form of
Proposition~\ref{prop:cartan-torus-regularity}. This basis makes
the first point canonical without changing the Cartan type.

Let $F=L(G[p^n])$ be the full splitting field, with Galois group
$\Gamma\subset C$, and let $x_F,y_F$ be the coordinates of
$v,Jv$. After choosing the special framing and a coefficient
map, Lemma~\ref{lem:scalar-compatible-frame} gives the equivariant
homomorphism $R_n(k_0)\to\mathcal O_F$ carrying $x,y$ to these
coordinates. All torsion points lie in the formal group because
the good special fibre is supersingular.

Extend $v_L$ to $F$ and put $q=c_0p^r/p^{2n-2}>0$.
Then $v_L(x_F)=qA_0$ and $v_L(y_F)=qB_0$. Since
$(1-p)A_0+pB_0=1$,
\[
 v_L(x_F^{1-p}y_F^p)=q>0,\qquad
 v_L(x_F^{B_0}/y_F^{A_0})=0.
\]
Thus $\pi_F=x_F^{1-p}y_F^p$ is integral of positive valuation
and $t_F=x_F^{B_0}/y_F^{A_0}$ is a unit. Substitution defines an
actual map
\[
 S_0=R_n(k_0)[\pi,t,t^{-1}]\longrightarrow\mathcal O_F.
\]
The images of $x,y$ are nonzero, so every rational expression used
here is defined. Its integrality follows from the displayed
valuations, without an assertion about fractional invariant lifts.
The residue of $t_F$ is some $c\in\kappa^\times$, and
$S_0/(\pi)=k_0[t,t^{-1}]$.

If $c$ is algebraic over $k_0$, then $c\in k_0$. Every denominator
outside $(\pi,t-[c])$ has nonzero residue, giving the actual local
map $B_c\to\mathcal O_F$ of the original closed chart.
If $c$ is transcendental over $k_0$, choose any $d\in k_0^\times$.
A denominator outside $(\pi,t-[d])$ reduces to a Laurent polynomial
$f(t)$ with $f(d)\ne0$. It is nonzero, so $f(c)\ne0$ as well.
Consequently there is still a ring map
$B_d\to\mathcal O_F$. In this second case the map is generally
not local at the closed point $d$; we do not apply a closed-point
locality assertion to it.

In either case, write $d=c$ in the first case and retain the chosen
$d$ in the second. The map on $B_d$ is $\Gamma$-equivariant.
Every $C$-invariant is $\Gamma$-invariant, even if $\Gamma$ is a
proper subgroup. Hence restriction gives a ring map
\[
                   A_d=B_d^C\longrightarrow
                   \mathcal O_F^\Gamma=\mathcal O_L.
\]
Use the exact identity inside the original invariant ring from
Proposition~\ref{prop:cartan-torus-regularity}:
\[
              p=u_d\rho_d^N,\qquad u_d\in A_d^\times,
              \qquad N=(p^4-1)p^{n-1}.
\]
A unit maps to a unit because its inverse maps into $\mathcal O_L$
as well. The image of $\rho_d$ is nonzero, since $p\ne0$, and has
positive integral valuation by the identity. Thus $N\mid e_L$
in both residue cases. This use of the unit identity requires no
locality at a transcendental closed point and no extension of
the chart theorem beyond $k_0$.

Finally,
\[
 e_L=(p^2-1)e=c_0(p^4-1)p^r,\qquad
 \frac{e_L}{N}=\frac{c_0}{p^2}
              =\frac{p-1}{2p^2}\in(0,1),
\]
which is impossible. The same argument applies to the second group.
\end{proof}

No equality between a point-orbit size and the full image order
was used. At $p=5,r=1$, the contradiction is
$e_L=6240<78000=N$. The proof excludes the nonsplit-Cartan image
hypothesis for an actual comparison; it does not force that
hypothesis for a general comparison.

\begin{proof}[Proof of Corollary~\ref{cor:boundary-nonsplit-self-twist}]
The completed unramified passage kills the unramified part of the
quadratic character. Its remaining tame part is killed by the
common extension of degree $p^2-1$. The intertwiner therefore
commutes with the image over $L$. By
Lemma~\ref{lem:quadratic-self-twist-cartan}, irreducible quadratic
reduction makes its centralizer the nonsplit quadratic \'etale
algebra. First torsion is rational over $L$, so the image is in
its principal-unit subgroup. This contradicts
Corollary~\ref{cor:boundary-nonsplit-cartan-image}.
\end{proof}

The split-intertwiner case is outside this deduction. The general
non-Cartan boundary configuration and its character index remain
open.
